\section{What was checked, and how}
\label{sec:scope}

This note reports a numerical audit of the Order-2 ``FK'' contribution to
the lensing two-point function, the term the draft plots in
\code{figures/analysis3\_NLO\_FFFK.pdf} and
\code{figures/analysis3\_cl\_full\_vs\_O0.pdf} and quotes in
\code{sections/insights.tex} and \code{sections/conclusion.tex}.

Six questions were asked, in this order.

\begin{enumerate}
\item Does the plotted FK curve evaluate the quantity the paper defines?
      (Section~\ref{sec:sampling})
\item Is the multipole cutoff in the vertex build converged, and if not,
      what does the converged answer rest on? (Section~\ref{sec:cutoff})
\item Does the line-of-sight floor $\lammin$ of the tabulated vertex bias
      the result, and does the bias depend on source redshift?
      (Section~\ref{sec:lambdamin})
\item Is the absolute normalisation of the vertex right, having never been
      compared against anything external? (Section~\ref{sec:normalisation})
\item Do the draft's statements about the E and B polarization structure
      survive a resolved vertex? (Section~\ref{sec:eb})
\item Does a nonlinear bispectrum give a more realistic prediction?
      (Section~\ref{sec:nonlinear}), and if the answer is no, what exactly
      is missing? (Section~\ref{sec:uv})
\end{enumerate}

The first is the one that changes published numbers. The fourth is the one
whose answer is unreservedly positive. The rest fall in between.

The first was settled by an independent reviewer who was given the figure
production path and nothing else. That reviewer was not told what any
earlier work had concluded, was blocked from reading the earlier notes, and
was told explicitly that finding nothing wrong would be a useful result. It
worked forward from the figure manifest through the generator scripts, the
folded \code{xi\_*.npz} files, the run configuration, the vertex callable
and the tabulated vertex, testing rather than assuming at each step. Its
conclusion agreed with the earlier work, and it added several results the
earlier work did not have. The comparison is recorded in
\code{notes/task\_a\_blind\_audit\_comparison.md}.

Two controls make the rest of the note interpretable.

\paragraph{The pipeline is reproducible.} Re-running the production FK fold
from a redirected copy of its own configuration reproduces the deployed
result bit for bit, with maximum relative difference zero across all six
component pairs. Every comparison below therefore isolates exactly one
change.

\paragraph{The rebuild machinery is exact.} The vertex table is rebuilt
here in pieces: one radial-and-multipole LOW branch, and a set of HIGH
multipole windows that are summed. This is only legitimate if the pieces
are independent, so that was tested rather than assumed. Table rows are
bit-exactly independent of one another in both branches: computing a row in
a table of six rows and in a table of two rows gives identical doubles
(\code{np.array\_equal} true, maximum relative difference $0.000$). The
multipole windows are disjoint in maximum leg, so their sum is the same
integral. Together these mean one build at the finest grid yields every
coarser grid by row selection, and every smaller cutoff by partial
summation, with no approximation. The reimplementation of the production
triple mesh reproduces the deployed table's 1671 rows exactly and in the
same order.

All commands are in \code{driver\_field\_emulators/code/}, and every number
quoted below names the script that produced it.
